%
% presence-heartbeat-blindput.tex  (FROZEN NEGATIVE CONTROL: do not edit)
%
% Z Specification for the heartbeat write of a presence row (PL-1704).
%
% Companion to nats-relay.tex: that specification models the sessions bucket
% as SessionStore (sessions : SESSIONKEY \pinj KVVALUE) and says nothing about
% two writers racing on one row.  This one models the race, and the
% compare-and-set that closes it.  It is kept apart because every component
% added to the composed state of nats-relay.tex multiplies a state space that
% already holds 831,363 states.
%
% Type-check: fuzz -t presence-heartbeat.tex
% Model-check: see the end of this file.
%

\documentclass[a4paper,10pt,fleqn]{article}
\usepackage[margin=1in]{geometry}
\usepackage{fuzz}
\usepackage[colorlinks=true,linkcolor=blue,citecolor=blue,urlcolor=blue]{hyperref}

\begin{document}

\title{Biff Presence Heartbeat: A Z Specification}
\author{Punt Labs}
\date{November 2026}
\maketitle

%%----------------------------------------------------------------------
\section{Introduction}
%%----------------------------------------------------------------------

A session's presence row lives in one key of the sessions bucket.  Two
writers use it.  The heartbeat (\texttt{NatsRelay.heartbeat}) reads the row,
moves \texttt{last\_active} forward, and writes the row back.  Everything
else (\texttt{plan}, \texttt{mesg}, a tty rename) reads the row, changes one
field, and writes it back.

The heartbeat's read and its write are two round trips.  A plan saved
between them is in the bucket when the heartbeat writes, and the heartbeat
writes the row it read: the plan is gone.  The loss is silent, and it is
the heartbeat's own tick that causes it.  PL-1704 is a report of exactly this.

The bucket gives every write a revision, and \texttt{kv.update} writes only
if the revision is the one the caller names.  The heartbeat now names the
revision it read.  A write that finds the row changed fails, and the
heartbeat reads again; after three failures it gives up until the next tick.
One skipped tick is harmless: the row lives for three days.

This specification states the property and checks it twice: against the
blind write that was shipped, where ProB finds the lost plan, and against
the compare-and-set, where the state is unreachable.

%%----------------------------------------------------------------------
\section{Basic Types}
%%----------------------------------------------------------------------

\begin{zed}
  ZBOOL ::= ztrue | zfalse
\end{zed}

\begin{zed}
  FIELD ::= fplan | fmesg
\end{zed}

$FIELD$ stands for every field of the row that a writer other than the
heartbeat owns.  Two of them are enough to say that the heartbeat preserves
\emph{every} such field and not one in particular.

\begin{axdef}
  maxRev : \nat \\
  maxBeat : \nat \\
  maxTries : \nat
  \where
  maxRev = 4 \\
  maxBeat = 2 \\
  maxTries = 3
\end{axdef}

$maxRev$ and $maxBeat$ bound the revision and the heartbeat clock so that
ProB can enumerate them.  $maxTries$ is the retry limit of
\texttt{heartbeat}.  Field values are the numbers $0$ and $1$.

\begin{zed}
  HbPhase ::= hbIdle | hbRead
\end{zed}

%%----------------------------------------------------------------------
\section{State}
%%----------------------------------------------------------------------

\begin{schema}{Row}
  fields : FIELD \fun (0 \upto 1) \\
  beat : 0 \upto maxBeat \\
  curRev : 0 \upto maxRev
\end{schema}

The row in the bucket: the fields the other writers own, the heartbeat clock
($last\_active$), and the revision the bucket assigned to the last write.

\begin{schema}{Heartbeat}
  phase : HbPhase \\
  seenFields : FIELD \fun (0 \upto 1) \\
  seenRev : 0 \upto maxRev \\
  tries : 0 \upto maxTries
\end{schema}

The heartbeat in flight.  After its read, $seenFields$ and $seenRev$ are the
row and the revision it holds; $tries$ counts the failed writes of this tick.

\begin{schema}{Intent}
  intended : FIELD \fun (0 \upto 1)
\end{schema}

A ghost.  $intended$ is what the other writers last wrote to each field, in
the order they wrote it.  It is not in the bucket and the code has no such
variable; it exists so that a lost write is a state and not an argument.

\begin{schema}{PresenceState}
  Row \\
  Heartbeat \\
  Intent
\end{schema}

\begin{schema}{PresenceInv}
  PresenceState
  \where
  fields = intended
\end{schema}

The property.  The row in the bucket holds, for every field, the last value
a writer other than the heartbeat gave it: the heartbeat never overwrites a
plan or any other field written after its read.

\begin{schema}{Invariant}
  PresenceInv
\end{schema}

\begin{schema}{Init}
  PresenceState'
  \where
  fields' = \{ fplan \mapsto 0, fmesg \mapsto 0 \} \\
  beat' = 0 \\
  curRev' = 0 \\
  phase' = hbIdle \\
  seenFields' = \{ fplan \mapsto 0, fmesg \mapsto 0 \} \\
  seenRev' = 0 \\
  tries' = 0 \\
  intended' = \{ fplan \mapsto 0, fmesg \mapsto 0 \}
\end{schema}

%%----------------------------------------------------------------------
\section{Operations}
%%----------------------------------------------------------------------

\begin{schema}{HeartbeatRead}
  \Delta Heartbeat \\
  \Xi Row \\
  \Xi Intent
  \where
  phase = hbIdle \\
  tries < maxTries \\
  phase' = hbRead \\
  seenFields' = fields \\
  seenRev' = curRev \\
  tries' = tries
\end{schema}

\texttt{kv.get} of the row.  Enabled while the tick has tries left.

\begin{schema}{WriteField}
  \Delta Row \\
  \Xi Heartbeat \\
  \Delta Intent \\
  f? : FIELD \\
  v? : 0 \upto 1
  \where
  curRev < maxRev \\
  fields' = fields \oplus \{ f? \mapsto v? \} \\
  intended' = intended \oplus \{ f? \mapsto v? \} \\
  beat' = beat \\
  curRev' = curRev + 1
\end{schema}

A writer other than the heartbeat sets one field.  It runs whenever it
likes, in particular while a heartbeat is in $hbRead$.

\begin{schema}{HeartbeatCommit}
  \Delta Row \\
  \Delta Heartbeat \\
  \Xi Intent
  \where
  phase = hbRead \\
  curRev < maxRev \\
  fields' = seenFields \\
  beat' = (beat + 1) \mod (maxBeat + 1) \\
  curRev' = curRev + 1 \\
  phase' = hbIdle \\
  seenFields' = seenFields \\
  seenRev' = seenRev \\
  tries' = 0
\end{schema}

\texttt{kv.update(key, row, last=entry.revision)}.  The guard $curRev =
seenRev$ is the compare: the bucket accepts the write only if no write has
landed since the read.  The row written is the row read with the clock
moved, $fields' = seenFields$; the guard is what makes that equal to the row
now in the bucket.  The clock wraps at $maxBeat$ only to keep ProB's state
space finite.

\begin{schema}{HeartbeatConflict}
  \Xi Row \\
  \Delta Heartbeat \\
  \Xi Intent
  \where
  phase = hbRead \\
  curRev \neq seenRev \\
  phase' = hbIdle \\
  seenFields' = seenFields \\
  seenRev' = seenRev \\
  tries' = tries + 1 \\
  tries < maxTries
\end{schema}

\texttt{KeyWrongLastSequenceError}: the row changed since the read.  Nothing
is written; the loop goes round again, and $HeartbeatRead$ is no longer
enabled once $tries = maxTries$ (``Session changed concurrently, skip
heartbeat'').

\begin{schema}{NextTick}
  \Xi Row \\
  \Delta Heartbeat \\
  \Xi Intent
  \where
  phase = hbIdle \\
  tries > 0 \\
  tries' = 0 \\
  phase' = hbIdle \\
  seenFields' = seenFields \\
  seenRev' = seenRev
\end{schema}

The next heartbeat tick starts with a fresh budget.

\begin{schema}{Quiesce}
  \Xi PresenceState
\end{schema}

The bound on $curRev$ ends the writers' history after four writes.  $Quiesce$
is the stutter that stands for ``nothing more happens''; without it ProB
reports the end of the bounded history as a deadlock, which says something
about the bound and nothing about the protocol.

%%----------------------------------------------------------------------
\section{Model Checking}
%%----------------------------------------------------------------------

\begin{enumerate}
  \item \textbf{No lost write}: $PresenceInv$, $fields = intended$, is a
    state invariant.  ProB checks it over the whole bounded state space.
  \item \textbf{Liveness of the tick}: $EF\; beat \neq 0$: a heartbeat
    can land.  A guard that were never enabled would make the first
    property vacuous.
  \item \textbf{The race is real}: $EF\;(phase = hbRead \land curRev \neq
    seenRev)$ is reachable, so the compare-and-set has something to refuse.
\end{enumerate}

The superseded write is reproduced in the frozen variant
\texttt{presence-heartbeat-blindput.tex}, whose $HeartbeatCommit$ has no
$curRev = seenRev$ guard (the \texttt{kv.put} the code used to make).  ProB
reports an invariant violation there, with the trace $HeartbeatRead \to
WriteField \to HeartbeatCommit$: the plan written between the read and the
write is overwritten by the value the heartbeat read.

%%----------------------------------------------------------------------
%% Reproducing the checks.  Run from the repo root; probcli is not assumed to
%% be on PATH.  Bounded run (one probcli at a time):
%%
%%   PROBCLI="${PROBCLI:-$HOME/Applications/ProB/probcli}"
%%   P="-p DEFAULT_SETSIZE 1 -p MAXINT 6 -p MAX_INITIALISATIONS 20 \
%%      -p MAX_OPERATIONS 500 -p TIME_OUT 600000"
%%   "$PROBCLI" docs/specs/presence-heartbeat.tex $P -model_check
%%   "$PROBCLI" docs/specs/presence-heartbeat.tex $P -model_check -ctlformula "EF {beat/=0}"
%%   "$PROBCLI" docs/specs/presence-heartbeat.tex $P -model_check -ctlformula \
%%       "EF {phase=hbRead & curRev/=seenRev}"
%%
%% Negative control, which must report an invariant violation:
%%
%%   "$PROBCLI" docs/specs/presence-heartbeat-blindput.tex $P -model_check
%%----------------------------------------------------------------------

\end{document}
